% ==============================
% Dummy Overleaf Project (HE) — Linear Algebra Book
% Exact environments: definition, example, exercise, remark,
% theorem, lemma, proposition, corollary, proof, solution, hint
% Numbering by chapter; solution/hint unnumbered.
% ==============================
\documentclass[12pt]{book}

%Comments (Nadav):
% Replace '||' with '\|' and better use \left\| \right\| where possible.
% There should be no '\L' in the file.
% Replace aleph, bet, ... with \begin{enumerate}[label=\alph*.]
% Replace 1,2,3,... with enumerate
% in general, use \left and \right as needed.
% usepackage{braket} \Set{ | } is your friend
% Adjust parentheses.
% Add shortcuts (newcommand) for \mathbb{R}, \vec{v}, \vec{w}, etc.
% Decide whether you're using $ math $ or \( math \)


% ---- Math & friends ----
\usepackage{amsmath,amssymb,mathtools}
\usepackage{amsthm}
\usepackage{hyperref}
\usepackage{graphicx}
\usepackage{booktabs}
\usepackage{enumitem}
\usepackage{todonotes}

% ---- Hebrew + fonts ----
\usepackage{fontspec}
\usepackage{polyglossia}
\setdefaultlanguage{hebrew}
\setotherlanguage{english}
% Overleaf has Frank Ruehl CLM in TeX Live
\newfontfamily\hebrewfont[Script=Hebrew]{Frank Ruehl CLM}
\newfontfamily\hebrewfontsf[Script=Hebrew]{Frank Ruehl CLM}
\newfontfamily\hebrewfonttt[Script=Hebrew]{Frank Ruehl CLM}


% ---- Theorem-like environments (by chapter) ----
\theoremstyle{plain}
\newtheorem{theorem}{משפט}[chapter]
\newtheorem{lemma}[theorem]{למה}
\newtheorem{proposition}[theorem]{טענה}
\newtheorem{corollary}[theorem]{מסקנה}

\theoremstyle{definition}
\newtheorem{definition}[theorem]{הגדרה}
\newtheorem{example}[theorem]{דוגמה}
\newtheorem{exercise}[theorem]{תרגיל}
\newtheorem{remark}[theorem]{הערה}

% Unnumbered for website folding
\newtheorem*{solution}{פתרון}
\newtheorem*{hint}{רמז}

% ---- Macros (optional) ----
\newcommand{\R}{\mathbb{R}}

\title{ספר אלגברה לינארית — דגם אוברליף}
\author{צוות הקורס}
\date{\today}

\begin{document}
\frontmatter
\maketitle
\tableofcontents
\mainmatter

\chapter{גיאומטריה במישור ובמרחב} 

בפרק זה נרצה להגדיר ולהבין וקטורים מבחינה גיאומטרית ואלגברית. נתמקד תחילה במישור (עם דגש על הישרים המוכלים בו) ואחר כך נעבור למרחב (שם נדבר גם על מישורים).

\subsection{ המישור \(\R^2\)}

בפרק \(0\) דיברנו על הישר הממשי וגם על המישור המרוכב. אפשר גם
לתאר את המישור (הממשי) בלי מספרים מרוכבים, אבל הרעיון דומה מאוד.

\begin{definition}
המישור \(\R^2\) הוא קבוצת כל הנקודות בעלות שתי
קוארדינטות ממשיות, כלומר
\begin{equation*}
	.\mathbb{R}^{2}=\{(x,y)|x,y\in\mathbb{R}\}
\end{equation*}
\end{definition}

\begin{example} \((0,0),(1,-2),(-3,-5),(\frac{3}{5},\sqrt{2}),(\pi,-10)\in\mathbb{R}^{2}\)

לעומת זאת: \(1,(i,0),(0,1,2)\notin\R^{2}\)
\end{example}
\begin{center}
	(מישור אינטראקטיבי שמאפשר ללחוץ על נקודה ולהציג את הוקטור עם הקוארדינטות
	שלו)
\end{center}

\begin{definition}
וקטור \(\vec{v}=(x,y)\) במישור הוא חץ שיוצא מהראשית \((0,0)\)
ומסתיים בנקודה \((x,y)\). 
\end{definition} 

\begin{remark}
נשתמש באותן הקוארדינטות לוקטור וגם לנקודה המתאימה. נבין את הכוונה
לפי ההקשר, אבל לרוב נדבר על וקטורים ולא על נקודות.
\end{remark}

הגדרנו וקטורים כחצים שיוצאים מהראשית, אבל הם יכולים לצאת מכל נקודה.
מה שקובע את הוקטור הם גודלו וכיוונו, לא נקודת ההתחלה. וקטורים זהים יופיעו כחצים מקבילים ושווי גודל במישור.
\begin{center}
	(מישור אינטראקטיבי שמאפשר להזיז וקטור נתון)
\end{center}

\begin{remark}
וקטורים שווים אם ורק אם יש שוויון בכל קוארדינטה, כלומר \((a,b)=(c,d)\)
אם ורק אם \(a=c\) וגם \(b=d\). למשל,
\((1,2)=(1,2)\)
אך \((1,2)\neq(1,-2)\). שני הוקטורים האחרונים שווים בגודלם אך הכיוונים שונים.
\end{remark}

\subsection{פעולות על וקטורים}

שתי הפעולות הבסיסיות שניתן לבצע על וקטורים הן חיבור וכפל בסקלר. סקלר
הוא מילה נרדפת למספר, וכאן נסמנו \(c\in\R\).

חיבור: \((a_{1},a_{2})+(b_{1},b_{2})=(a_{1}+b_{1},a_{2}+b_{2})\)

כפל בסקלר: \(c\cdot(a_{1},a_{2})=(ca_{1},ca_{2})\) 

\begin{remark}
הוקטור הנגדי (שווה גודל אך הפוך בכיוון) מתקבל ע"י
כפל בסקלר \(-1\). חיסור בין שני וקטורים זה כמו חיבור בין הראשון
לגרסה הנגדית של השני, כלומר:
\begin{equation*}
	(a_{1},a_{2})-(b_{1},b_{2})=(a_{1},a_{2})+(-1)\cdot(b_{1},b_{2})=(a_{1}-b_{1},a_{2}-b_{2})
\end{equation*}
\end{remark}

\begin{example}\ 
\begin{enumerate}[label=\alph*.]
    \item \((1,2)+(10,20)=(1+10,2+20)=(11,22)\)
    \item \(2\cdot(5,3)=(10,6)\)
    \item \(-5\cdot(2,5)=(-10,-25)\)
    \item \((9,2)-(11,8)=(9-11,2-8)=(-2,-6)\)
\end{enumerate}
\end{example}

אפשר להבין חיבור וקטורי באופן גיאומטרי ע"י כלל המקבילית.
כלל זה קובע שאם נצייר שני וקטורים \(\vec{v},\vec{w}\) שיוצאים מהראשית
ונשלים אותם למקבילית ע"י הוספת שני וקטורים מקבילים
ושווי גודל כצלעות נגדיות, אז האלכסון שיוצא מהראשית יתאר את \(\vec{v}+\vec{w}\).
ההפרש \(\vec{v}-\vec{w}\) מיוצג ע"י האלכסון השני. 
\begin{center}
	(מישור אינטראקטיבי שמראה חיבור/הפרש וקטורי כאלכסוני מקבילית)
\end{center}

זה גם מוביל אותנו לנוסחה שמתארת וקטור שמתחיל בנקודה \((a_{1},a_{2})\)
ומסתיים ב-\((b_{1},b_{2})\): מקבלים את ההפרש הוקטורי \((b_{1}-a_{1},b_{2}-a_{2})\).

\begin{exercise} 
הראו שהוקטור שמתחיל ב-\((1,3)\) ומסתיים ב-\((5,0)\), שווה לוקטור שמתחיל ב-
\((5,-1)\)
ומסתיים ב-
\((9,-4)\).
\end{exercise}

כפל בסקלר זו פעולה שאפשר לפרש גיאומטרית כמתיחה של הוקטור אם \(c>1\),
או כיווצו אם \(0<c<1\), בלי לשנות כיוון. כאשר הסקלר שלילי, יש היפוך
כיוון ובנוסף מתיחה/כיווץ בהתאם לערך של \(|c|\).
\begin{center}
(מישור אינטראקטיבי שמראה כפל בסקלר כמתיחה/כיווץ/היפוך כיוון)
\end{center}

\subsection{גדלים וזוויות ב-\(\R^{2}\)}

כאמור, לכל וקטור יש גודל וכיוון. אם נניח שהוא מתחיל בראשית, אז גודלו/אורכו
הוא המרחק מנקודת הסיום לראשית. הכיוון נקבע לפי הזווית \(\theta\)
עם הכיוון החיובי של ציר \(x\) בדיוק כמו בקוארדינטות פולריות.
עבור וקטור \((a,b)\) נקבל את הנוסחאות הבאות:

גודל: \(||(a,b)||=\sqrt{a^{2}+b^{2}}\)

זווית: \(\tan\theta=\frac{b}{a}\)

כבר דיברנו על החישוב המדויק של \(\theta\) בסוף הפרק הקודם.
לא ממש נצטרך לחשב זוויות בקורס, אבל זה טוב לידע כללי.

\begin{remark}
כל וקטור נקבע ביחידות ע"י הגודל והכיוון שלו. לכן
קוארדינטות פולריות יכולות להחליף קוארדינטות קרטזיות, אבל בקורס שלנו
יותר נוח להשתמש בקוארדינטות קרטזיות וכך נעשה. לקראת סוף הקורס נחזור לגיאומטריה והחישוב של גדלים יהיה חשוב.
\end{remark}

\begin{example}
א.
\(||(2,1)||=\sqrt{2^{2}+1^{2}}=\sqrt{5}\)

ב. 
\(||(3,-4)||=\sqrt{3^{2}+(-4)^{2}}=5\)

ג. 
\(||(-3,-7)||=\sqrt{(-3)^{2}+(-7)^{2}}=\sqrt{58}\)
\end{example}

אם כפל בסקלר מתאים למתיחה/כיווץ, זה אמור להתבטא בגודל. מכאן הטענה
הבאה:

\begin{proposition}\label{proposition:norm-scalar}
\(||c\cdot(a,b)||=|c|\cdot||(a,b)||\) לכל וקטור \((a,b)\in\R^{2}\)
ולכל סקלר \(c\in\R\).
\end{proposition}

\begin{proof}
נשתמש בהגדרות ונחשב:
\begin{align*}
||c\cdot(a,b)||&=||(ca,cb)||=\sqrt{(ca)^{2}+(cb)^{2}}=\sqrt{c^{2}(a^{2}+b^{2})}=\sqrt{c^2}\sqrt{a^2+b^2}\\
&=|c|\sqrt{a^{2}+b^{2}}=|c|\cdot||(a,b)||
\end{align*}
נציין שמתקיים \(\sqrt{c^2}=|c|\) כי מדובר בשורש חיובי.
\end{proof}

\begin{remark}
בפרט, כפל ב-\(-1\) לא משנה את גודל הוקטור. רק הכיוון מתהפך.
\end{remark}

\subsection{ישרים ב-\(\R^{2}\)}

נרצה להבין את הקשר בין ישרים לוקטורים. יותר קל להבין את זה במקרה של
ישרים שעוברים דרך הראשית. ניקח ישר כזה ונבחר נקודה כלשהי שאינה הראשית,
מה שייתן לנו וקטור. הישר הוא אוסף כל הנקודות שמתקבלות ע"י
כפל בסקלר של וקטור זה.

\begin{example}
הישר שמשוואתו \(y=2x\) עובר דרך הראשית והנקודה 
\((1,2)\).
דרך נוספת לתאר את הישר היא כקבוצת כל הכפולות (בסקלר ממשי) של הוקטור
\((1,2)\). כל הוקטורים האלה מתאימים לנקודות על הישר, וכותבים את הנקודות כמו הוקטורים (כי הם יוצאים מהראשית). כך נקבל את ההצגה הבאה לישר כקבוצת נקודות:
	\begin{equation*}
	\{t(1,2)|t\in\mathbb{R}\}=\{(t,2t)|t\in\mathbb{R}\}
	\end{equation*}
\end{example}

\begin{center}
	(אנימציה שמראה איך נקודות הישר מתקבלות ע"י כפל בסקלר
	של הוקטור)
\end{center}

אכן, זה מיידי שכל נקודה מהצורה \((t,2t)\) מקיימת את המשוואה \(y=2x\). 

באופן כללי, הישר שעובר דרך הראשית בכיוון הוקטור \(\vec{v}\) נתון
ע"י ההצגה \(\{t\vec{v}|t\in\R\}\). נשים לב
שעבור \(t<0\) מקבלים וקטורים עם כיוון מנוגד (הצד השני של הישר).

נכליל: ישר שעובר דרך \((x_{0},y_{0})\) בכיוון הוקטור \(\vec{v}\)
נתון ע"י ההצגה
\begin{equation*}
	 .\{(x_{0},y_{0})+t\vec{v}|t\in\mathbb{R}\}
\end{equation*}
אם נתון שהישר גם עובר דרך \((x_{1},y_{1})\), אז נקבל
\begin{equation*}
	.\{(x_{0},y_{0})+t(x_{1}-x_{0},y_{1}-y_{0})|t\in\mathbb{R}\}
\end{equation*}

\begin{exercise}
איזו נקודה מתקבלת עבור \(t=0\)? \(t=1\)? \(t=2\)?
\end{exercise}

\begin{remark}
ההצגה לעיל נקראת הצגה פרמטרית של ישר. לעומת זאת, הצגה מהצורה
\begin{equation*}
    L=\{(x,y)\in\mathbb{R}^{2}|ax+by=c\}
\end{equation*}
נקראת הצגה אלגברית של
ישר. הצגה אלגברית אינה יחידה, כי למשל המשוואה \(y=2x\) שקולה למשוואה
\(2y=4x\). באופן דומה, גם הצגה פרמטרית אינה יחידה משתי סיבות: אפשר
לבחור את נקודת ההתחלה \((x_{0},y_{0})\) באופן שרירותי כל עוד היא
על הישר, וגם אפשר לבחור את \((x_{1},y_{1})\) כרצוננו ובהתאם את
וקטור הכיוון \(\vec{v}\) (אפשר לקחת כל כפולה שלו בסקלר שאינו 
\(0\)).
כך מקבלים אינסוף הצגות פרמטריות שונות לאותו הישר. למשל:
\begin{equation*}
	\{t(1,2)|t\in\mathbb{R}\}=\{(1,2)+s(2,4)|s\in\mathbb{R}\}=\{(2,4)+r(-3,-6)|r\in\mathbb{R}\}
\end{equation*}
\end{remark}

\begin{exercise}
עבור הישר שמשוואתו \(y=3x+1\), מצאו שלוש הצגות פרמטריות
שונות. נסו לשנות גם את נקודת ההתחלה וגם את וקטור הכיוון.
\end{exercise}

\subsection{מכפלה סקלרית}

מכפלה סקלרית היא פעולה נוספת בין וקטורים, אבל התוצאה שלה היא סקלר
ומכאן שמה.

\begin{definition}\label{def:scalar-product-plane}
עבור וקטורים \((a_1,a_2),(b_1,b_2)\in\R^2\)
נגדיר \((a_{1},a_{2})\cdot(b_{1},b_{2})=a_{1}b_{1}+a_{2}b_{2}\).
\end{definition}

\begin{example}
א. 
\((1,3)\cdot(2,5)=1\cdot2+3\cdot5=17\)

ב. 
\((3,-2)\cdot(6,5)=3\cdot6+(-2)\cdot5=8\)

ג. 
\((1,5)\cdot(-5,1)=1\cdot(-5)+5\cdot1=0\)
\end{example}


למכפלה סקלרית יש קשר לגודל. יש מספר תכונות שכדאי לזכור:

\begin{proposition}
יהיו \(\vec{v_{1}}=(x_{1},y_{1}),\,\vec{v_{2}}=(x_{2},y_{2}),\,\vec{v_{3}}=(x_{3},y_{3})\in\R^{2}\)
וקטורים. אז מתקיים:

א. 
\(\vec{v_{1}}\cdot\vec{v_{2}}=\vec{v_{2}}\cdot\vec{v_{1}}\).

ב. 
\((\vec{v_{1}}+\vec{v_{3}})\cdot\vec{v_{2}}=\vec{v_{1}}\cdot\vec{v_{2}}+\vec{v_{3}}\cdot\vec{v_{2}}\).

ג. \(\vec{v_{1}}\cdot\vec{v_{1}}=||\vec{v_{1}}||^{2}\).

ד. \(\vec{v_{1}}\cdot(c\vec{v_{2}})=c(\vec{v_{1}}\cdot\vec{v_{2}})\) לכל סקלר \(c\in\R\).

ה. \(||\vec{v_{1}}||=0\iff\vec{v_{1}}=(0,0)\).

ו. \(||\frac{\vec{v_{1}}}{||\vec{v_{1}}||}||=1\) אם \(\vec{v_{1}}\neq(0,0)\).
\end{proposition}

\begin{proof}
נסתפק בסעיפים ג' ו-ה'. סעיפים אחרים יופיעו כתרגילים. לפי ההגדרה של מכפלה סקלרית, מתקיים
\begin{equation*}
.\vec{v_{1}}\cdot\vec{v_{1}}=x_{1}^{2}+y_{1}^{2}=\left(\sqrt{x_{1}^{2}+y_{1}^{2}}\right)^{2}=||\vec{v_{1}}||^{2}
\end{equation*}

סכום ריבועים (של מספרים ממשיים) הוא אי-שלילי, והוא מתאפס אם ורק אם
כל מחובר מתאפס. במקרה זה \(x_{1}^{2}=y_{1}^{2}=0\) גורר \(x_{1}=y_{1}=0\),
או באופן שקול \(\vec{v_{1}}=(0,0)\).
\end{proof}

מכפלה סקלרית גם קשורה לזווית בין שני וקטורים. ליתר דיוק, הכוונה היא
לזווית הקטנה ביניהם (בין 
\(0\) ל-\(\pi\) ברדיאנים).

\begin{center}
	(מישור אינטראקטיבי שמאפשר ליצור שני וקטורים ולקבל את המכפלה הסקלרית
	וגם את הזווית ברדיאנים)
\end{center}

\begin{proposition}
יהיו \(\vec{v},\vec{w}\) שני וקטורים ב-\(\mathbb{R}^{2}\)
השונים מ-
\((0,0)\), ותהי \(\alpha\) הזווית הקטנה ביניהם. אז
מתקיים
\begin{equation}\label{eq:cosine-scalar}
	.\cos\alpha=\frac{\vec{v}\cdot\vec{w}}{||\vec{v}|||\cdot||\vec{w}||}
\end{equation}
\end{proposition}
ההוכחה תדגים היטב את השימוש בתכונות של מכפלה סקלרית.

\begin{proof}
נסתמך על משפט הקוסינוסים מטריגונומטריה. נסתכל על המשולש שנוצר
ע"י שני הוקטורים:
\begin{center}
	(ציור עם משולש בעל צלעות שמתאימות לשני הוקטורים וגם לוקטור ההפרש מול
	\(\alpha\))
\end{center}

\begin{equation*}
||\vec{v}||^{2}+||\vec{w}||^{2}=||\vec{v}-\vec{w}||^{2}+2||\vec{v}||\cdot||\vec{w}||\cos\alpha
\end{equation*}

נשתמש בתכונות של מכפלה סקלרית כדי למצוא את הקשר בין \(||v-w||^{2}\)
ל-\(v\cdot w\):
\begin{equation*}
	||\vec{v}-\vec{w}||^{2}=(\vec{v}-\vec{w})\cdot(\vec{v}-\vec{w})=\vec{v}\cdot\vec{v}-\vec{v}\cdot\vec{w}-\vec{w}\cdot\vec{v}+\vec{w}\cdot\vec{w}=||\vec{v}||^{2}-2\vec{v}\cdot\vec{w}+||\vec{w}||^{2}
\end{equation*}

נציב את המשוואה השנייה בראשונה ונקבל
\begin{equation*}
.||\vec{v}||^{2}+||\vec{w}||^{2}=||\vec{v}||^{2}-2\vec{v}\cdot\vec{w}+||\vec{w}||^{2}+2||\vec{v}||\cdot||\vec{w}||\cos\alpha
\end{equation*}

נקזז \(||\vec{v}||^{2}+||\vec{w}||^{2}\) מכל אגף ונקבל
\begin{equation*}
.0=-2\vec{v}\cdot\vec{w}+2||\vec{v}||\cdot||\vec{w}||\cos\alpha\implies\cos\alpha=\frac{\vec{v}\cdot\vec{w}}{||\vec{v}||\cdot||\vec{w}||}
\end{equation*}

כנדרש.
\end{proof}

\begin{corollary}
	וקטורים \(\vec{v},\vec{w}\) השונים מ-\((0,0)\) הם מאונכים/ניצבים
	אם ורק אם \(\vec{v}\cdot\vec{w}=0\). במקרה זה נסמן \(\vec{v}\perp\vec{w}\).
\end{corollary}

\begin{proof}
תהי \(\alpha\) הזווית הקטנה ביניהם. אז מתקיים
\begin{equation*}
	.\vec{v}\cdot\vec{w}=0\iff\cos\alpha=\frac{\vec{v}\cdot\vec{w}}{||\vec{v}||\cdot||\vec{w}|}=0\iff\alpha=\frac{\pi}{2}
\end{equation*}
\end{proof}

\begin{example}
א. 
\((1,2)\perp(-2,1)\) כי \((1,2)\cdot(-2,1)=0\).

ב.
באופן כללי, לכל \((a,b)\in\R^{2}\) מתקיים \((a,b)\perp(-b,a)\)
כי \((a,b)\cdot(-b,a)=0\).
\end{example}

\begin{remark}
הוקטור \((0,0)\) הוא מקרה חריג כי הוא חסר כיוון. אבל מאחר
שמתקיים \(\vec{v}\cdot(0,0)=0\) לכל וקטור \(\vec{v}\in\R^2\),
עדיין נגדיר \(\vec{v}\perp(0,0)\). זהו הוקטור היחיד שמאונך לכל
וקטור במישור, כי כל וקטור אחר לא מאונך לעצמו.
\end{remark}

\subsection{מעבר בין הצגה פרמטרית של ישר להצגה אלגברית}

ההצגה האלגברית של ישר קשורה באופן הדוק למכפלה סקלרית. נבין את הקשר
דרך דוגמה.

\begin{example}
נסתכל על הישר \(L=\{t(2,-1)|t\in\mathbb{R}\}\). רואים לפי
ההצגה הפרמטרית שוקטור הכיוון הוא \((2,-1)\). לפי המכפלה הסקלרית
אפשר לקבוע שהוקטור \((1,2)\) מאונך לו. הוא לא רק מאונך ל-\((2,-1)\)
אלא לכל וקטור שמתאים לנקודה על הישר, כי לכל \(t\in\R\)
מתקיים
\begin{equation*}
	.(1,2)\cdot(2t,-t)=2t-2t=0
\end{equation*}
\end{example}

\begin{center}
	(הישר עם וקטור הנורמל)
\end{center}

להיפך, אפשר לתאר את הישר כאוסף כל הוקטורים שמאונכים לוקטור \((1,2)\).
כלומר כל וקטור \((x,y)\in L\) נתון ע"י המשוואה
\begin{equation*}
	.(x,y)\cdot(1,2)=0\iff x+2y=0
\end{equation*}

או בכתיב קבוצות: \(L=\{(x,y)\in\mathbb{R}^{2}|x+2y=0\}\).

\begin{definition}
וקטור המאונך לכל וקטור שמתאים לשתי נקודות על ישר (התחלה וסיום)
נקרא נורמל לישר.
\end{definition}

\begin{center}
(מישור אינטראקטיבי שמאפשר לסובב ישר כאשר הנורמל מסתובב יחד איתו בהתאם)
\end{center}

באופן כללי, לישר עם הצגה פרמטרית \(\{(x_{0},y_{0})+t(a,b)|t\in\R\}\)
יש וקטור נורמל \((-b,a)\) המאונך לוקטור הכיוון \((a,b)\). אבל
וקטור הנורמל אינו יחיד, כי אם נכפיל אותו בסקלר \(c\neq0\) הוא יישאר
מאונך לוקטור הכיוון של הישר. כל שני וקטורים מאונכים יישארו מאונכים
גם אם נכפיל את אחד מהם (או שניהם) בסקלר.

כדי לעבור להצגה אלגברית נשתמש בוקטור שמתחיל ב-\((x_{0},y_{0})\)
ומסתיים ב-\((x,y)\), כלומר הוקטור \((x-x_{0},y-y_{0})\). וקטור
זה בהכרח מאונך לוקטור הנורמל \((-b,a)\), ולכן צריך להתקיים:
\begin{align*}
	(x-x_{0},y-y_{0})\cdot(-b,a)=0&\iff-b(x-x_{0})+a(y-y_{0})=0\\
    &\iff-bx+ay=-bx_{0}+ay_{0}
\end{align*}
כלומר קיבלנו הצגה אלגברית \(L=\{(x,y)\in\mathbb{R}^{2}|-bx+ay=-bx_{0}+ay_{0}\}\).
נדגיש שהבחירה של נקודת ההתחלה \((x_{0},y_{0})\) היא שרירותית.
המשוואה מראה שהביטוי \(-bx+ay\) קבוע לאורך הישר, וקבוע זה מופיע
באגף ימין של המשוואה.

\begin{example}
נסתכל על הישר \(L=\{(2,3)+t(1,-1)|t\in\mathbb{R}\}\). יש
לו וקטור נורמל \((1,1)\) ולכל \((x,y)\in L\) הוקטור \((x-2,y-3)\)
מאונך ל-\((1,1)\), ולכן מתקיים
\begin{equation*}
	.(x-2,y-3)\cdot(1,1)=0\iff x-2+y-3=0\iff x+y=5
\end{equation*}

קיבלנו הצגה אלגברית \(L=\{(x,y)\in\R^2|x+y=5\}\).
\end{example}

\begin{exercise}
מצאו הצגה אלגברית של \(L=\{(1,-2)+t(3,1)|t\in\R\}\).
\end{exercise}

\subsection{מעבר בין הצגה אלגברית של ישר להצגה פרמטרית}

כדי למצוא הצגה פרמטרית של ישר, צריך למצוא שתי נקודות עליו. האחת בשביל
נקודת ההתחלה, והשנייה בשביל נקודת הסיום של וקטור הכיוון. הבחירה של
נקודות אלו היא שרירותית לחלוטין, אבל יחסית נוח להשתמש בנקודות חיתוך
עם הצירים.

\begin{example}
נסתכל על הישר \(L=\{(x,y)\in\mathbb{R}^{2}|2x+3y=1\}\).
אם נציב \(y=0\) נקבל \(x=\frac{1}{2}\), ולכן נבחר את \((\frac{1}{2},0)\)
כנקודת התחלה. אם נציב \(x=0\) נקבל \(y=\frac{1}{3}\), ולכן נבחר
את \((0,\frac{1}{3})\) כנקודת סיום של וקטור הכיוון. כך נקבל וקטור
כיוון \((0,\frac{1}{3})-(\frac{1}{2},0)=(-\frac{1}{2},\frac{1}{3})\).

מכאן \(L=\{(\frac{1}{2},0)+t(-\frac{1}{2},\frac{1}{3})|t\in\mathbb{R}\}\)
היא הצגה פרמטרית אחת מתוך רבות. מי שלא אוהב שברים יכול להכפיל ב-\(6\)
את וקטור הכיוון, וגם אפשר להחליף את נקודת ההתחלה ב-\((2,-1)\),
למשל. אז גם מתקיים \L{$L=\{(2,-1)+t(-3,2)|t\in\mathbb{R}\}$} ושתי
התשובות נכונות באותה המידה. נשים לב שלפי הנוסחה \L{$(-b,a)$}, וקטור
הנורמל שמתאים להצגה הפרמטרית הראשונה הוא \L{$(-\frac{1}{3},-\frac{1}{2})$}.
באופן דומה, וקטור הנורמל שמתאים להצגה הפרמטרית השנייה הוא \L{$(-2,-3)$}.
שני הוקטורים האלה הם כפולות של הוקטור \L{$(2,3)$} שמופיע בתוך ההצגה
האלגברית המקורית בתור המקדמים של \L{$x,y$} בהתאמה.
\end{example}

באופן כללי, עבור הצגה אלגברית \L{$L=\{(x,y)\in\mathbb{R}^{2}|ax+by=c\}$}
נבחר נקודת התחלה \L{$(x_{0},y_{0})$} כלשהי. זו נקודה שמקיימת את משוואת
הישר ולכן בהכרח \L{$ax_{0}+by_{0}=c$}, ולאחר הצבה במשוואה נקבל\L{%
	\[
	.(x,y)\in L\iff ax+by=ax_{0}+by_{0}\iff a(x-x_{0})+b(y-y_{0})=0\iff(x-x_{0},y-y_{0})\perp(a,b)
	\]
}%

אז \L{$(a,b)$} הוא וקטור נורמל של \L{$L$}. כל וקטור כיוון שנמצא
)ע\textquotedbl י בחירה נקודה נוספת על \L{$L$}( בהכרח מאונך לוקטור
זה. זו בדיוק המשמעות של התנאי \L{$(x-x_{0},y-y_{0})\perp(a,b)$}.

\begin{exercise}
מצאו הצגה פרמטרית של \L{$L=\{(x,y)\in\mathbb{R}^{2}|3x+y=5\}$}.
האם וקטור הכיוון מאונך ל-\L{$(3,1)$}?
\end{exercise}


\section{הקדמה}
דוגמה לפרק שנכתב לפי ההנחיות. הקפידו להשתמש בשם הסביבה המדויק ולהוסיף תיוגים עם קידומת
כמו \verb|def:|, \verb|thm:|, \verb|exr:|, \verb|eq:| וכן הלאה.

\begin{definition}\label{def:vectorspace}
מרחב וקטורי מעל שדה \(F\) הוא \dots
\end{definition}

\begin{theorem}\label{thm:basis-exists}
לכל מרחב וקטורי סופי-ממד מעל \(F\) קיים בסיס.
\end{theorem}
\begin{proof}[הוכחה]
רעיון ההוכחה: נבנה קבוצה בת הוספה עד שאינה ניתנת להגדלה, ואז נראה שהיא בסיס. \qedhere
\end{proof}

\begin{example}\label{ex:first}
ב-\(\R^2\) הקבוצה \(\{(1,0),(0,1)\}\) היא בסיס סטנדרטי.
\end{example}

\begin{exercise}\label{exr:find-basis}
מצאו בסיס ל-\(\R^2\) המורכב מהוקטורים \((1,1)\) ו-\((1,-1)\).
\end{exercise}

\begin{hint}
אפשר להתחיל בבדיקה של אי-תלות לינארית בעזרת דטרמיננטה.
\end{hint}

\begin{solution}
הדטרמיננטה של המטריצה \(\begin{pmatrix}1&1\\ 1&-1\end{pmatrix}\) שווה \(-2\neq 0\), ולכן הוקטורים בת\"ל.
מכאן שהם גם פורשים את \(\R^2\), ולכן מהווים בסיס כנדרש.
\end{solution}

ראו משפט~\ref{thm:basis-exists}, הגדרה~\ref{def:vectorspace}, ודוגמה~\ref{ex:first}.

\section{משוואות}
נשתמש ב-\verb|equation| ו-\verb|align| (לא \verb|eqnarray|).
\begin{equation}\label{eq:idempotent}
P^2 = P
\end{equation}
\begin{align}\label{eq:two-steps}
Ax &= b\\
x  &= A^{-1}b
\end{align}
הפניה למשוואות: \eqref{eq:idempotent}, \eqref{eq:two-steps}.

\backmatter
\end{document}